% ============================================================
% tikz preamble.tex
% Shared TikZ colors, sizes, and figure macros for the whole paper.
% Loaded once from main.tex. Individual figures live in the
% 'fig - ...' files and only use what is defined here.
% ============================================================

% ---------- shared colors ----------
\definecolor{Zcol}{HTML}{FF1744}
\definecolor{Xcol}{HTML}{00B7FF}
\definecolor{HCol}{HTML}{D35400}
\definecolor{SwapCol}{HTML}{8E44AD}
\colorlet{Rcol}{Zcol}
\colorlet{Bcol}{Xcol}
% luminance-balanced stabilizer tints (annotated diagrams, Figure 3)
\colorlet{Xtint}{Xcol!26}
\colorlet{Ztint}{Zcol!17}
% luminance-balanced vivid fills (lattice maps, Figure 2)
\colorlet{Xvivid}{Xcol!95}
\colorlet{Zvivid}{Zcol!50}
\newcommand{\qdLine}{0.7pt}  % qudit circle rim weight

% ---------- shared sizes ----------
\newcommand{\qdR}{4.8pt}      % qudit circle radius
\newcommand{\numFont}{\scriptsize} % numbers inside circles
\newcommand{\labelFont}{\normalsize} % X/Z labels
\newcommand{\arrAngle}{70}    % wedge arrow half-angle
\newcommand{\arrLen}{0.17cm}  % wedge arrow tail length
\newcommand{\edge}{1.20cm}
\newcommand{\thick}{1pt}

% ---------- five-qutrit patch and stabilizer macros (d=2 figure) ----------
% Outlined stabilizer-patch fills used in the d=2 redefined X/Z grids.
% These match the d=3 convention: each stabilizer is shown as a lightly
% filled patch with its own colored perimeter, while X- and Z-type checks
% are drawn on separate copies of the lattice.
\newcommand{\FiveQBlueTriOutlined}[3]{%
  \path[fill=Xtint,draw=Xcol,line width=1.0pt]
    (#1) -- (#2) -- (#3) -- cycle;
}

\newcommand{\FiveQRedTriOutlined}[3]{%
  \path[fill=Ztint,draw=Zcol,line width=1.0pt]
    (#1) -- (#2) -- (#3) -- cycle;
}

% On the d=2 redefined grids, one stabilizer boundary would otherwise run
% straight through qudit 2 even though qudit 2 is not in that stabilizer.
% These helper paths replace that diagonal by a clean detour around the
% visible q2 silhouette. The blue detour keeps the circular line-around-node
% shape from the previous version. The red detour follows the left outline of
% the q2 arrow/circle silhouette so B_1 reads as bypassing q2 rather than using
% q2.
\newcommand{\FiveQDetourR}{\qdR}
\newcommand{\FiveQArrowTipDetour}{\dimexpr\qdR+\arrLen\relax}
\newcommand{\FiveQDetouredBlueTriQoneQthreeQfour}{%
  % A_2 bypasses q2. Fill the detoured patch, draw the normal boundary, then
  % overdraw only the q2-contacting arc underneath q2 so the visible blue sits
  % flush against the black circle with no white halo.
  \path[fill=Xtint,draw=none,line cap=round,line join=round]
    (q1)
    -- ($(q2)+(45:\FiveQDetourR)$)
    arc[start angle=45,end angle=-135,radius=\FiveQDetourR]
    -- (q3) -- (q4) -- cycle;
  \draw[Xcol,line width=1.0pt,line cap=round,line join=round]
    (q1)
    -- ($(q2)+(45:\FiveQDetourR)$)
    arc[start angle=45,end angle=-135,radius=\FiveQDetourR]
    -- (q3) -- (q4) -- cycle;
  \draw[Xcol,line width=1.8pt,line cap=round,line join=round]
    ($(q2)+(45:\FiveQDetourR)$)
    arc[start angle=45,end angle=-135,radius=\FiveQDetourR];
}
\newcommand{\FiveQDetouredRedTriQzeroQoneQthree}{%
  % B_1 bypasses q2. The contact path exactly follows the q2 silhouette:
  % right wedge base -> wedge tip -> left wedge base -> left-lower circle arc.
  % That contact path is overdrawn slightly wider underneath q2; after q2 is
  % redrawn, the visible red has zero white gap to the black qudit/arrow.
  \path[fill=Ztint,draw=none,line cap=round,line join=round]
    (q0) -- (q1)
    -- ($(q2)+(20:\FiveQDetourR)$)
    -- ($(q2)+(90:\FiveQArrowTipDetour)$)
    -- ($(q2)+(160:\FiveQDetourR)$)
    arc[start angle=160,end angle=225,radius=\FiveQDetourR]
    -- (q3) -- cycle;
  \draw[Zcol,line width=1.0pt,line cap=round,line join=round]
    (q0) -- (q1)
    -- ($(q2)+(20:\FiveQDetourR)$)
    -- ($(q2)+(90:\FiveQArrowTipDetour)$)
    -- ($(q2)+(160:\FiveQDetourR)$)
    arc[start angle=160,end angle=225,radius=\FiveQDetourR]
    -- (q3) -- cycle;
  \draw[Zcol,line width=1.8pt,line cap=round,line join=round]
    ($(q2)+(20:\FiveQDetourR)$)
    -- ($(q2)+(90:\FiveQArrowTipDetour)$)
    -- ($(q2)+(160:\FiveQDetourR)$)
    arc[start angle=160,end angle=225,radius=\FiveQDetourR];
}

% === AUTO-GENERATED routed stabilizer macros (d=2) ===
%% AUTO-GENERATED d=2 routed stabilizers
\newcommand{\FiveQRouteX}{%
  \path[fill=Xtint,draw=none] (0.0000,0.0000)--(1.0000,1.0000)--(-1.0000,1.0000)--cycle;
  \path[fill=Xtint,draw=none] (-1.0000,-1.0000)--(1.0000,-1.0000)--(1.0000,1.0000)--(0.1423,0.1423)--(0.1492,0.1309)--(0.1559,0.1197)--(0.1625,0.1085)--(0.1688,0.0975)--(0.1751,0.0863)--(0.1812,0.0751)--(0.1868,0.0634)--(0.1905,0.0510)--(0.1934,0.0385)--(0.1955,0.0257)--(0.1968,0.0129)--(0.1972,0.0000)--(0.1968,-0.0129)--(0.1955,-0.0257)--(0.1934,-0.0385)--(0.1905,-0.0510)--(0.1868,-0.0634)--(0.1822,-0.0755)--(0.1769,-0.0872)--(0.1708,-0.0986)--(0.1640,-0.1096)--(0.1565,-0.1201)--(0.1483,-0.1300)--(0.1395,-0.1395)--(0.1300,-0.1483)--(0.1201,-0.1565)--(0.1096,-0.1640)--(0.0986,-0.1708)--(0.0872,-0.1769)--(0.0755,-0.1822)--(0.0634,-0.1868)--(0.0510,-0.1905)--(0.0385,-0.1934)--(0.0257,-0.1955)--(0.0129,-0.1968)--(0.0000,-0.1972)--(-0.0129,-0.1968)--(-0.0257,-0.1955)--(-0.0385,-0.1934)--(-0.0510,-0.1905)--(-0.0634,-0.1868)--(-0.0755,-0.1822)--(-0.0872,-0.1769)--(-0.0986,-0.1708)--(-0.1096,-0.1640)--(-0.1201,-0.1565)--(-0.1300,-0.1483)--(-0.1395,-0.1395)--cycle;
  \draw[Xcol,line width=1.0pt,line cap=round,line join=round] (0.0000,0.0000) -- (1.0000,1.0000) -- (-1.0000,1.0000) -- cycle;
  \draw[Xcol,line width=1.0pt,line cap=round,line join=round] (-1.0000,-1.0000) -- (1.0000,-1.0000) -- (1.0000,1.0000) -- (0.1423,0.1423) -- (0.1492,0.1309) -- (0.1559,0.1197) -- (0.1625,0.1085) -- (0.1688,0.0975) -- (0.1751,0.0863) -- (0.1812,0.0751) -- (0.1868,0.0634) -- (0.1905,0.0510) -- (0.1934,0.0385) -- (0.1955,0.0257) -- (0.1968,0.0129) -- (0.1972,0.0000) -- (0.1968,-0.0129) -- (0.1955,-0.0257) -- (0.1934,-0.0385) -- (0.1905,-0.0510) -- (0.1868,-0.0634) -- (0.1822,-0.0755) -- (0.1769,-0.0872) -- (0.1708,-0.0986) -- (0.1640,-0.1096) -- (0.1565,-0.1201) -- (0.1483,-0.1300) -- (0.1395,-0.1395) -- (0.1300,-0.1483) -- (0.1201,-0.1565) -- (0.1096,-0.1640) -- (0.0986,-0.1708) -- (0.0872,-0.1769) -- (0.0755,-0.1822) -- (0.0634,-0.1868) -- (0.0510,-0.1905) -- (0.0385,-0.1934) -- (0.0257,-0.1955) -- (0.0129,-0.1968) -- (0.0000,-0.1972) -- (-0.0129,-0.1968) -- (-0.0257,-0.1955) -- (-0.0385,-0.1934) -- (-0.0510,-0.1905) -- (-0.0634,-0.1868) -- (-0.0755,-0.1822) -- (-0.0872,-0.1769) -- (-0.0986,-0.1708) -- (-0.1096,-0.1640) -- (-0.1201,-0.1565) -- (-0.1300,-0.1483) -- (-0.1395,-0.1395) -- cycle;
}
\newcommand{\FiveQRouteZ}{%
  \path[fill=Ztint,draw=none] (-1.0000,-1.0000)--(-0.1395,-0.1395)--(-0.1483,-0.1300)--(-0.1565,-0.1201)--(-0.1640,-0.1096)--(-0.1708,-0.0986)--(-0.1769,-0.0872)--(-0.1822,-0.0755)--(-0.1868,-0.0634)--(-0.1905,-0.0510)--(-0.1934,-0.0385)--(-0.1955,-0.0257)--(-0.1968,-0.0129)--(-0.1972,-0.0000)--(-0.1968,0.0129)--(-0.1955,0.0257)--(-0.1934,0.0385)--(-0.1905,0.0510)--(-0.1868,0.0634)--(-0.1812,0.0751)--(-0.1751,0.0863)--(-0.1688,0.0975)--(-0.1625,0.1085)--(-0.1559,0.1197)--(-0.1492,0.1309)--(-0.1423,0.1423)--(-0.1351,0.1540)--(-0.1275,0.1661)--(-0.1194,0.1787)--(-0.1109,0.1921)--(-0.1017,0.2063)--(-0.0918,0.2216)--(-0.0809,0.2385)--(-0.0689,0.2572)--(-0.0554,0.2784)--(-0.0399,0.3029)--(-0.0218,0.3319)--(-0.0000,0.3672)--(0.0218,0.3319)--(0.0399,0.3029)--(0.0554,0.2784)--(0.0689,0.2572)--(0.0809,0.2385)--(0.0918,0.2216)--(0.1017,0.2063)--(0.1109,0.1921)--(0.1194,0.1787)--(0.1275,0.1661)--(0.1351,0.1540)--(0.1423,0.1423)--(1.0000,1.0000)--(-1.0000,1.0000)--cycle;
  \path[fill=Ztint,draw=none] (1.0000,-1.0000)--(1.0000,1.0000)--(0.0000,0.0000)--cycle;
  \draw[Zcol,line width=1.0pt,line cap=round,line join=round] (-1.0000,-1.0000) -- (-0.1395,-0.1395) -- (-0.1483,-0.1300) -- (-0.1565,-0.1201) -- (-0.1640,-0.1096) -- (-0.1708,-0.0986) -- (-0.1769,-0.0872) -- (-0.1822,-0.0755) -- (-0.1868,-0.0634) -- (-0.1905,-0.0510) -- (-0.1934,-0.0385) -- (-0.1955,-0.0257) -- (-0.1968,-0.0129) -- (-0.1972,-0.0000) -- (-0.1968,0.0129) -- (-0.1955,0.0257) -- (-0.1934,0.0385) -- (-0.1905,0.0510) -- (-0.1868,0.0634) -- (-0.1812,0.0751) -- (-0.1751,0.0863) -- (-0.1688,0.0975) -- (-0.1625,0.1085) -- (-0.1559,0.1197) -- (-0.1492,0.1309) -- (-0.1423,0.1423) -- (-0.1351,0.1540) -- (-0.1275,0.1661) -- (-0.1194,0.1787) -- (-0.1109,0.1921) -- (-0.1017,0.2063) -- (-0.0918,0.2216) -- (-0.0809,0.2385) -- (-0.0689,0.2572) -- (-0.0554,0.2784) -- (-0.0399,0.3029) -- (-0.0218,0.3319) -- (-0.0000,0.3672) -- (0.0218,0.3319) -- (0.0399,0.3029) -- (0.0554,0.2784) -- (0.0689,0.2572) -- (0.0809,0.2385) -- (0.0918,0.2216) -- (0.1017,0.2063) -- (0.1109,0.1921) -- (0.1194,0.1787) -- (0.1275,0.1661) -- (0.1351,0.1540) -- (0.1423,0.1423) -- (1.0000,1.0000) -- (-1.0000,1.0000) -- cycle;
  \draw[Zcol,line width=1.0pt,line cap=round,line join=round] (1.0000,-1.0000) -- (1.0000,1.0000) -- (0.0000,0.0000) -- cycle;
}
% === end routed macros ===

\newcommand{\FiveQPatchRedefinedXGrid}{%
  \def\L{1.0}
  \coordinate (q0) at (-\L, \L); \coordinate (q1) at ( \L, \L);
  \coordinate (q2) at ( 0 , 0 );
  \coordinate (q3) at (-\L,-\L); \coordinate (q4) at ( \L,-\L);
  \FiveQRouteX
  \FiveQPatchEdges
}

\newcommand{\FiveQPatchRedefinedZGrid}{%
  \def\L{1.0}
  \coordinate (q0) at (-\L, \L); \coordinate (q1) at ( \L, \L);
  \coordinate (q2) at ( 0 , 0 );
  \coordinate (q3) at (-\L,-\L); \coordinate (q4) at ( \L,-\L);
  \FiveQRouteZ
  \FiveQPatchEdges
}

% standard black wedge arrow
\newcommand{\WedgeArrow}[2]{%
  \path[fill=black,draw=none]
    ($(#1)+({\qdR*cos(#2+\arrAngle)},{\qdR*sin(#2+\arrAngle)})$)
      arc[start angle=#2+\arrAngle, end angle=#2-\arrAngle, radius=\qdR]
      -- ($(#1)+({(\qdR+\arrLen)*cos(#2)},{(\qdR+\arrLen)*sin(#2)})$) -- cycle;
}

% ----- PATCH GEOMETRY -----

% 1. Edges Only
\newcommand{\FiveQPatchEdges}{%
  \def\L{1.0} \def\thick{1.0pt}
  \coordinate (q0) at (-\L, \L); \coordinate (q1) at ( \L, \L);
  \coordinate (q2) at ( 0 , 0 );
  \coordinate (q3) at (-\L,-\L); \coordinate (q4) at ( \L,-\L);
  \draw[line width=\thick] (q0) -- (q1);
  \draw[line width=\thick] (q3) -- (q4);
  \draw[line width=\thick] ( 0,\L) -- (0,-\L);
}

% 2. Filled (Reference)
\newcommand{\FiveQPatchFilled}{%
  \def\L{1.0} \def\thick{1.0pt}
  \coordinate (q0) at (-\L, \L); \coordinate (q1) at ( \L, \L);
  \coordinate (q2) at ( 0 , 0 );
  \coordinate (q3) at (-\L,-\L); \coordinate (q4) at ( \L,-\L);
  \path[fill=Xvivid,draw=none] (q0) -- (q2) -- (q1) -- cycle;
  \path[fill=Xvivid,draw=none] (q3) -- (q2) -- (q4) -- cycle;
  \path[fill=Zvivid,draw=none] (q0) -- (q2) -- (q3) -- cycle;
  \path[fill=Zvivid,draw=none] (q1) -- (q2) -- (q4) -- cycle;
  \draw[line width=\thick] (q0) -- (q1);
  \draw[line width=\thick] (q3) -- (q4);
  \draw[line width=\thick] ( 0,\L) -- (0,-\L);
}

% NEW: 2b. Overlapped Fill
\newcommand{\FiveQPatchOverlapped}{%
  \def\L{1.0} \def\thick{1.0pt}
  \coordinate (q0) at (-\L, \L); \coordinate (q1) at ( \L, \L);
  \coordinate (q2) at ( 0 , 0 );
  \coordinate (q3) at (-\L,-\L); \coordinate (q4) at ( \L,-\L);
  % Fill specific triangles with Blue
  \path[fill=Bcol,opacity=0.55,draw=none] (q0) -- (q2) -- (q3) -- cycle;
  \path[fill=Bcol,opacity=0.55,draw=none] (q2) -- (q3) -- (q4) -- cycle;
  % Fill specific triangles with Red
  \path[fill=Rcol,opacity=0.55,draw=none] (q0) -- (q2) -- (q3) -- cycle;
  \path[fill=Rcol,opacity=0.55,draw=none] (q2) -- (q3) -- (q4) -- cycle;
  % Draw lines
  \draw[line width=\thick] (q0) -- (q1);
  \draw[line width=\thick] (q3) -- (q4);
  \draw[line width=\thick] ( 0,\L) -- (0,-\L);
}

% 3. Circles + Wedges
\newcommand{\FiveQPatchCircles}{%
  \foreach \P/\txt in {q0/0,q1/1,q2/2,q3/3,q4/4}{
    \draw[line width=\qdLine,fill=white] (\P) circle (\qdR);
    \node[font=\numFont] at (\P) {\txt};
  }
  \WedgeArrow{q0}{0}\WedgeArrow{q1}{0}\WedgeArrow{q3}{0}\WedgeArrow{q4}{0}\WedgeArrow{q2}{90}
}

% 3b. Circles + Wedges (Smaller Font for Part d, e, f)
\newcommand{\FiveQPatchCirclesSmall}{%
  \foreach \P/\txt in {q0/0,q1/1,q2/2,q3/3,q4/4}{
    \draw[line width=\qdLine,fill=white] (\P) circle (\qdR);
    \node[font=\numFont] at (\P) {\txt};
  }
  \WedgeArrow{q0}{0}\WedgeArrow{q1}{0}\WedgeArrow{q3}{0}\WedgeArrow{q4}{0}\WedgeArrow{q2}{90}
}

% 4. Circles + Wedges Layered
\newcommand{\FiveQPatchCirclesOneUnderThree}{%
  \foreach \P/\txt in {q0/0,q2/2,q4/4}{
    \draw[line width=\qdLine,fill=white] (\P) circle (\qdR);
    \node[font=\numFont] at (\P) {\txt};
  }
  \def\off{0.10}
  \coordinate (qoneBR) at ($(q3)+(\off,-\off)$);
  \draw[line width=\qdLine,fill=white] (qoneBR) circle (\qdR);
  \node[font=\numFont] at (qoneBR) {1};
  \draw[line width=\qdLine,fill=white] (q3) circle (\qdR);
  \node[font=\numFont] at (q3) {3};
  \WedgeArrow{q0}{0}\WedgeArrow{q3}{0}\WedgeArrow{q4}{0}\WedgeArrow{q2}{90}\WedgeArrow{qoneBR}{0}
}
